\section{Introduction}
\label{sec:intro}

A flexible job shop does not run on machines alone. Jobs must be moved between
machines, and in an automated shop that movement is performed by a fleet of
automated guided vehicles whose availability, routing, and load state determine
whether a machine that is ready to work has anything to work on. Deciding the
two layers separately, fixing the production plan first and dispatching vehicles
to execute it afterwards, is the conventional decomposition. It is exact only when
transport is never the limiting resource. In an automated shop it can be, and
the coupling is then exactly what the decomposition discards: the production
plan determines which transport tasks exist and when they are due, and the
dispatch decisions determine whether the machines are ever fed.

The integrated problem is known in the recent literature as the
production--logistics collaborative scheduling problem
(PLCSP)~\citep{li2025plcsp,shi2025nested}, and in an older line as integrated
scheduling of machines and AGVs~\citep{ulusoy1993}. Both lines are active
today. Reinforcement learning is the dominant recent tool, applied through a
variety of representations: multi-agent PPO over hand-built priority
vectors~\citep{li2025plcsp}, nested hierarchical
architectures~\citep{shi2025nested}, heterogeneous graph
attention~\citep{taagnet2026,yang2026hgamppo,yue2027hchgnn}, and hybrid
GNN--Transformer encoders~\citep{gao2026gnntransformer}. A parallel line
attacks the same coupling with hand-designed search such as large-neighbourhood
search and Benders decomposition, and solves closely related three-layer
problems well~\citep{chen2025carseq,chen2026benders}. The two lines differ in
what they cost to extend: the search-based methods are tailored per problem,
and a change to the constraints generally means redesigning the encoding or the
cuts, whereas a learned policy is in principle indifferent to which constraints
are switched on.

That second property is the reason this paper exists, and it is also why the
field's current state is worth examining closely. Across the recent work, the
architecture has converged while the constraint coverage has not. Congestion
and path conflict are excluded in at least eleven papers, three of which name
the cause rather than the choice: an encoder that carries only connectivity
cannot represent congestion, so a policy using it cannot learn to avoid
it~\citep{gao2026gnntransformer,taagnet2026,li2025plcsp}. Preventive
maintenance is listed as future work rather than modelled~\citep{zhang2026amr}.
Batching as a decision appears in warehouse order picking but not in transport
inside a job shop~\citep{wang2025gphh}. Charging is modelled in several
works~\citep{huang2025collaborative,dong2026handling}, but whether the
constraint ever activates under the parameters used is rarely reported.

Two consequences follow, and they set up the contribution. First, a model can
carry a constraint that never fires. A constraint that is modelled but never
violated, or modelled but never actionable, contributes nothing to the problem
and nothing to the comparison. Whether it fires is reported rarely, one
constraint at a time, and never --- in the work we surveyed --- for every
modelled constraint of a benchmark family at once; the closest case counts the
actions one constraint triggers, on a single
instance~\citep{huang2025collaborative}.
Section~\ref{sec:results:active} reports activation for all ten constraints
across the whole benchmark. Second, an architecture
claim and a constraint-handling claim are different claims, and the second is
where the field is uneven. Section~\ref{sec:related} documents both.

Figure~\ref{fig:motivation} shows the decomposition failing on a three-machine
shop, and what changes when the two layers are decided together.

\begin{figure*}[htbp]
\centering
\includegraphics[width=\textwidth]{fig_motivation.pdf}
\caption{Deciding production and transport separately breaks down when the
vehicle fleet is the limiting resource; deciding them in one chain does not.
The shop has three machines and one vehicle. In (a) the machine schedule is
fixed first, with the vehicle absent from its input: it places two transfers
one time unit apart, job $J_1$ from $M_1$ to $M_2$ and job $J_2$ from $M_2$ to
$M_3$, because it assumes that transport is always available. In (b) the
vehicle is dispatched to execute that plan, and one vehicle cannot serve both
transfers: $J_1$ waits while the vehicle carries $J_2$, $M_2$ has nothing to
work on, and every later block on that machine shifts, so the loss appears as
idle machine time rather than as a worse machine sequence. In (c) both layers
are decided in one chain, so the plan is built knowing what the fleet can
deliver, and a single tour serves $J_1$ and then $J_2$ without idling any
machine. The example is constructed to make the mechanism visible and is not
drawn from a measured run.}
\label{fig:motivation}
\end{figure*}

This paper therefore spends its effort on the constraints and on the
representation they need. The policy's parts are standard: a Transformer
encoder, a candidate-scoring head per decision kind, and group-relative policy
optimization. What the paper adds to them is a formulation and a
representation. The formulation is the joint chain, which carries every
production and logistics decision of an episode under one log-probability and
one advantage. The representation is the transport-network-aware encoding of
Section~\ref{sec:method:state}, which puts the aisle network into the state
rather than leaving it to be inferred from the machine and vehicle tokens.
Every constraint then gets a mechanism that gives the policy something to act
on: a decision point where a rule used to decide, a dynamics change where the
model was incomplete, or an input feature where the state was not visible.
Every switch defaults to off with bit-identical behaviour, so that an ablation
measures the mechanism and not an incidental change elsewhere.
Section~\ref{sec:related:coverage} sets out, for each of the ten constraints,
what prior work does and what this paper does.

\paragraph{What we find}

On the fully constrained problem the learned policy dominates every
out-of-sample point of an NSGA-II front: the front's best re-evaluated values
are 79.56 for makespan, 6.239 for energy, and 20.73 for tardiness, against
78.90, 6.160, and 5.042 for the policy, and all twelve front points are
dominated. On the unconstrained problem the ordering reverses and NSGA-II edges
ahead, 49.14 against 49.27, on a problem where it has five thousand
deterministic evaluations available. That pair of results locates the method's
advantage in its representation of the constrained state rather than in its
search behaviour, which is the claim the rest of the paper examines.

Two further observations are worth stating before the results. An ablation over
the advantage shows that subtracting a group-mean baseline does measurable work
while dividing by the group standard deviation does not, which is a narrower
statement than the one usually made about group-relative advantages. And a
hypothesis we set out to confirm (that ablating on the makespan alone
systematically misses mechanisms) did not survive its own test. One property of
the reported objectives is not an observation at all and is stated where the
results are given: the energy account cannot move independently of the
makespan, because most of it is tied to the horizon by definition
(Section~\ref{sec:results:energy}).

\paragraph{Contributions}

\begin{enumerate}
\item \textbf{A joint-chain formulation, and a measurement of which part of it
      matters.} Production and logistics decisions in one episode are treated
      as one decision chain with one log-probability and one group-relative
      advantage, rather than as separate agents. We give the reason (group
      relative advantage compares samples within a group, and a group drawn
      from two policies is not a group) and then measure it. The
      measured answer is narrower than the usual statement: subtracting the
      group mean matters, dividing by the group standard deviation does not
      (Section~\ref{sec:results:advantage}). The arrangement is not merely a
      convenience: a recent comparison of joint and modular training on a job
      shop with transport resources finds a coordination gap, in which the
      modular arrangement does not recover what joint training
      obtains~\citep{link2026coordination}.

\item \textbf{A transport-network-aware encoding, and a mechanism for each of
      ten constraints.} The encoders in this literature carry the network's
      structure but not its live state: three papers state that theirs cannot
      represent congestion, so a policy using one cannot learn to avoid it.
      Where the network does enter an encoder it enters as static topology. We
      therefore make the transport network itself part of the state. The
      aisle segments become tokens carrying their own occupancy, a
      distance-dependent bias enters the attention score, and the route head
      reads the segments each candidate traverses. The aisle stops
      being something the policy is told about and becomes something it can act
      on. The remaining mechanisms follow the same discipline: a decision point
      where a rule used to decide, a dynamics change where the model was
      incomplete, or an input feature where the state was not visible.
      Section~\ref{sec:related:coverage} sets out the ten cases one by one and
      separates the ones where the difference is categorical from the ones where
      it is a matter of degree. Every switch defaults to off and reproduces the
      pre-mechanism trajectory bit for bit, which is what makes the ablations
      measurements rather than comparisons against an unstated baseline. We do
      not claim that the encoding transfers to shop layouts it was not trained
      on; the mechanisms are demonstrated on the layout of the instance.

\item \textbf{Results reported against interest.} We report the ablation table
      as a grouping rather than as an ordering, because the measured
      training-seed standard deviation is the same order as most of the
      differences in it. We report the hypothesis about single-objective
      ablation that did not hold, and the constraints that do not activate
      under the reported parameters, so that the coverage claim is a statement
      about the model rather than about demonstrated effects.
\end{enumerate}

\paragraph{Scope} The evaluation in this paper is a first pass. It uses one
instance of the MKT benchmark and a single training run per configuration, with
five seeds per arm reserved for the advantage ablation because that comparison
is about the training procedure itself. Section~\ref{sec:discussion} states
what this bounds. The experimental protocol is designed so that widening it
later changes the tables rather than the claims.

The remainder of the paper is organized as follows.
Section~\ref{sec:related} reviews the two literature lines, sets out the
constraint-by-constraint coverage, and explains why published figures from this
literature are not reproduced in our tables. Section~\ref{sec:problem}
formulates the problem, the ten constraints, and the three objectives.
Section~\ref{sec:method} describes the policy, the training procedure, and the
mechanism for each constraint. Section~\ref{sec:setup} gives the benchmark, the
configuration, and the reproducibility discipline.
Section~\ref{sec:results} reports the results, including the one negative one.
Section~\ref{sec:discussion} states the limitations, and
Section~\ref{sec:conclusion} concludes.
